% Last Updated: October 1, 2023
	\documentclass[a4paper, twoside]{article}
	% Margins
	\usepackage{geometry}
		\geometry{
			% showframe,
			left   = 20mm,
			right  = 20mm,
			top    = 20mm,
			bottom = 20mm
		}

	% Headers and footers
	\usepackage{fancyhdr}
		\pagestyle{fancy}
		\fancyhf{}
		\fancyhead[LE, RO]{\thepage}
		\fancyhead[CE]{\textsc{Cong Wen}}
		\fancyhead[CO]{\textsc\rightmark}

		\setlength{\headheight}{15pt}
		\renewcommand\headrulewidth{0pt}








\input{/Users/wenc/Documents/Math/universal-pre}

\title{\tbf{Mathematical physics}}
\author{\tbf{Notes} by Cong Wen, 02/2024}
\date{}





\begin{document}
\maketitle
\thispagestyle{empty}
\begin{abstract}
	This is my self-use notes\footnote{Last updated on February 23, 2024} in mathematical physics.

	Please reach out if you find anything incorrect.

\end{abstract}
\tableofcontents
% ----------------------------------------------------------------------------------------------------

\begin{dfn}[Dan's note, Axiom System 1.1]
	A physical system consists the data of
	\begin{enr}
		\item (States) A real convex set $S$. Extreme points of $S$ are \tbf{pure states}, others are \tbf{mixed states}.
		\item (Observables) A complex topological vector space $\cO$ with
			\begin{enr}
				\item an antilinear involution defining a real structure $\cO_{\bR}$,
				\item a dense subspace $\cO^{\ift}$ with a complex Lie algebra structure on $\cO^{\ift}$
				\item a map $\mathrm{Borel}(\bR,\bR)\tms\cO_{\bR}\ar\cO_{\bR}$.
			\end{enr}
			Elements of $\cO_{\bR}$ are called \tbf{observables}.
		\item (Measurement) A map $\cO_{\bR}\tms S\ar \mathrm{Prob}(\bR)$, sending $(A,\sgm)$ to a probability measure $\sgm_{A}$ on $\bR$.
		\item (Evolution) Maps from $\cO_{\bR}^{\ift}$ to one-parameter subgroups of $\Aut(S)$ and $\Aut(\cO)$.
	\end{enr}
	which satisfyies compatibility requirements.
\end{dfn}


% ----------------------------------------------------------------------------------------------------
\bibliographystyle{amsalpha}
\bibliography{/Users/wenc/Documents/Math/universal-ref}
\end{document}

